\documentclass[11pt,a4paper]{article}
% Common preamble for the hint sheets (\input from each file)
\usepackage[utf8]{inputenc}
\usepackage[T1]{fontenc}
\usepackage{lmodern}
\usepackage{amsmath,amssymb,amsthm,mathtools}
\usepackage[a4paper,margin=2.3cm]{geometry}
\usepackage{enumitem}
\usepackage{xcolor}
\usepackage[hidelinks]{hyperref}
\renewcommand{\leq}{\leqslant}
\renewcommand{\le}{\leqslant}
\renewcommand{\geq}{\geqslant}
\renewcommand{\ge}{\geqslant}
\setlength{\parindent}{0pt}
\setlength{\parskip}{0.35em}
% \hintstitle{Sheet title}{Class work / Homework}
\newcommand{\hintstitle}[2]{\begin{center}{\LARGE\bfseries Hints: #1\par}\vspace{0.5em}{\large #2\par}\end{center}
\noindent\textit{The hints to each problem are cumulative. Read one, try again, and open the next only if you are still stuck.}\par\medskip\hrule\medskip}
\newcommand{\problem}[1]{\par\bigskip{\large\bfseries Problem #1}\par\nopagebreak\smallskip}
\newcommand{\subpart}[1]{\par\smallskip\textbf{(#1)}\par\nopagebreak}
\newcommand{\hint}[1]{\par\smallskip\noindent\textcolor{black!70}{\textbf{Hint #1.}}\enspace}
\newcommand{\note}[1]{\par\smallskip\noindent\textit{Note.} #1\par}
\providecommand{\lt}{<}
\providecommand{\gt}{>}
\providecommand{\tightlist}{\setlength{\itemsep}{0pt}\setlength{\parskip}{0pt}}
\emergencystretch=3em

\begin{document}
\hintstitle{Probability -- 2. Geometric probability and random points}{Class work}

\problem{1}
\hint{1} Parametrise a line by its unit normal direction $\theta\in[0,\pi)$ and signed distance $p\in\mathbb R$ from the origin, with measure $dp\,d\theta$ (invariant under rigid motions). For a segment of length $\ell$ the lines in direction $\theta$ that meet it form a $p$-interval of length $\ell|\cos(\theta-\theta_0)|$; integrate to get $2\ell$.
\hint{2} For a convex polygon, almost every line meeting it crosses its boundary in exactly two points. Summing over the edges, $\sum_{\text{edges}}\mu(\text{lines meeting the edge})=2\mu(\text{lines meeting }K)$, hence $\mu=\text{perimeter}$; pass to general convex bodies by approximation. On the other hand $\mu=\int_0^\pi w(\theta)\,d\theta$, where $w(\theta)$ is the length of the projection on the normal direction, so the mean projection length is $\text{perimeter}/\pi$ (Cauchy).
\hint{3} Buffon: lines at spacing $d$, needle of length $\ell\le d$. The needle crosses a line iff its projection on the normal direction covers a point of $d\mathbb Z$, which, given the direction, has probability (projection length)$/d$. Averaging, $\Pr=\frac{2\ell}{\pi d}$ (the mean projection of a segment is $\frac{2\ell}{\pi}$, by the previous hint with perimeter $2\ell$).
\hint{4} In $\mathbb R^n$ the projection of a segment of length $\ell$ on a random line has length $\ell|u_1|$, $u$ uniform on $S^{n-1}$. Compute $\mathrm E|u_1|$ from the density of $u_1$, proportional to $(1-t^2)^{(n-3)/2}$: you get $\mathrm E|u_1|=\frac{\Gamma(n/2)}{\sqrt\pi\,\Gamma((n+1)/2)}$ (check: $2/\pi$ for $n=2$, $1/2$ for $n=3$).

\problem{2}
\hint{1} Write both probabilities with exchangeability. With $n=2019$ and $T=\triangle A_1A_2A_3$, $\Pr(\text{triangle})=\binom n3\mathrm E\big[(|T|/\pi)^{n-3}\big]$, and $\Pr(\text{quadrilateral})=\binom n4\Pr(Q)$, where $Q$ is the event that $A_1,\dots,A_4$ are in convex position and all other points lie inside their hull. Note $\binom n4/\binom n3=\frac{n-3}4=504$.
\hint{2} Condition on $A_1,A_2,A_3$. Let $R(T)$ be the set of positions of $A_4$ for which $A_1A_2A_3A_4$ is a convex quadrilateral (the three regions beyond a side and inside the angle at the opposite vertex). Then the quadrilateral contains $T$, so $\Pr(\text{quad})\ge\binom n4\mathrm E\big[\frac{|R(T)\cap D|}{\pi}(|T|/\pi)^{n-4}\big]$. It suffices that $504\,|R(T)\cap D|\ge|T|$, except on negligible configurations.
\hint{3} Let $BC=a$ be the longest side, $M$ its midpoint. Since $|OB|,|OC|\le1$, $|OM|^2\le1-\frac{a^2}4$, so the distance from $M$ to the circle is at least $1-|OM|\ge\frac{a^2}8$. Put $Y$ on ray $AM$ at distance $\frac{a^2}8$ beyond $M$: the triangle $BCY$ lies in $R(T)\cap D$ and $|BCY|/|T|=MY/MA\ge\frac{a^2/8}{(\sqrt3/2)a}=\frac{a}{4\sqrt3}$.
\hint{4} So the needed pointwise inequality holds whenever $a\ge4\sqrt3/504$; triangles with smaller longest side have area $<10^{-4}$ and contribute at most $(10^{-4})^{2016}$ to $\mathrm E|T|^{2016}$, which is negligible against the contribution of triangles of area $\ge1$ (positive probability), where the pointwise inequality holds with a large margin. Conclusion: the quadrilateral is more probable.

\problem{3}
\hint{1} A point of the circle is never inside a polygon inscribed in the same circle, so every vertex of $F$ is a crossing point of a red edge and a blue edge. Look at the cyclic sequence of colours.
\hint{2} Let the colours form $j$ red runs and $j$ blue runs around the circle. A red edge joining consecutive red points crosses the blue polygon iff a blue run lies between its endpoints, and then (if $j\ge2$) exactly twice. Hence $m=2j=$ number of colour changes if $j\ge2$, while for $j=1$ a chord separates the two hulls and $m=0$.
\hint{3} By linearity, $\mathrm E(\#\text{colour changes})=(n+k)\cdot\frac{2nk}{(n+k)(n+k-1)}=\frac{2nk}{n+k-1}$, and $\Pr(j=1)=\frac{n+k}{\binom{n+k}k}$. Therefore $\mathrm E\,m=\frac{2nk}{n+k-1}-\frac{2\,n!\,k!}{(n+k-1)!}$.

\problem{4}
\hint{1} Let $\alpha_1,\dots,\alpha_5$ be the arcs $AB,BC,CD,DE,EA$. Given $A$, the other four points are i.i.d.\ uniform, so $(\alpha_1,\dots,\alpha_5)$ is uniform on the simplex $\sum\alpha_i=2\pi$ and hence exchangeable.
\hint{2} Both triangles have base $AC$. Express their base angles by inscribed angles: in $ACB$ they are $\alpha_2/2$ (at $A$) and $\alpha_1/2$ (at $C$); in $ACK$ they are $\angle DAC=\alpha_3/2$ and $\angle ECA=\alpha_5/2$. A triangle on base $AC$ with base angles $x,y$ has height $AC/(\cot x+\cot y)$.
\hint{3} So the event is $\cot\frac{\alpha_1}2+\cot\frac{\alpha_2}2<\cot\frac{\alpha_3}2+\cot\frac{\alpha_5}2$. Swapping $(\alpha_1,\alpha_2)\leftrightarrow(\alpha_3,\alpha_5)$ preserves the distribution and turns the event into its complement (up to a null set). The probability is $\frac12$.

\problem{5}
\hint{1} Fix $P$. Since $\Delta$ lies inside $\Omega$, the chord $PQ$ meets $\Delta$ iff the ray $PQ$ does, i.e.\ iff $Q$ lies in the arc cut out by the two tangents from $P$ to $\Delta$. By the inscribed angle theorem this arc has angular size $4\arcsin\frac1{|PC|}$, $C=(r,0)$.
\hint{2} Hence the probability is $\frac2\pi\cdot\frac1{2\pi}\int_\Omega\arcsin\frac1{|P-C|}\,d\theta$, the circle average of $u(x)=\arcsin\frac1{|x|}$ evaluated at $x=P-C$. Study this average as a function of the centre $C$.
\hint{3} Check that $u$ is subharmonic for $|x|>1$: for $g(\rho)=\arcsin(1/\rho)$ one gets $g''+g'/\rho=(\rho^2-1)^{-3/2}>0$. The circle average of a subharmonic function, as a function of $C$, is subharmonic and radial, and a smooth radial subharmonic function $v(r)$ is non-decreasing ($rv'$ is non-decreasing and vanishes at $0$). The minimum is at $r=0$.

\end{document}
